\subsection{Клиентское веб-приложение (Frontend)}

Серверный API, описанный выше, используется клиентским веб-приложением, обеспечивающим пользовательский интерфейс системы. Клиентское приложение реализовано на SvelteKit 2 с Svelte 5~\cite{Svelte5} и представляет собой одностраничное приложение (\acrshort{spa}) с серверным рендерингом (\acrshort{ssr}). На рисунке~\ref{frontend-arch} представлена архитектура клиентского приложения.

\begin{figure}
  \centering
  \includegraphics[width=0.95\textwidth,height=0.85\textheight,keepaspectratio]{frontend-structure.png}
  \caption{Архитектура клиентского веб-приложения}
  \label{frontend-arch}
\end{figure}

\subsubsection{Маршрутизация и навигация}

Приложение содержит более 24 страниц, сгруппированных по функциональным областям: дашборд, животные, надои, воспроизводство, кормление, фитнес, пастьба, отчёты, аналитика, настройки, контакты, местоположения, быки, переводы, запасы, оповещения, задачи, ветеринария, финансы, рентабельность, поиск.

Навигация реализована через боковую панель (Sidebar) с иконками из библиотеки Lucide и группировкой разделов. Адаптивная верстка на основе Tailwind CSS обеспечивает корректное отображение на мобильных устройствах: на экранах менее 768px боковая панель сворачивается в гамбургер-меню.

\subsubsection{Аутентификация на стороне клиента}

Серверный хук (hooks.server.ts) проверяет наличие \acrshort{jwt} в cookie при каждом запросе. Декодируется полезная нагрузка токена и проверяется срок действия. Неаутентифицированные пользователи перенаправляются на страницу входа, аутентифицированные --- не могут перейти на страницу входа.

Управление состоянием аутентификации реализовано через Svelte-стор (auth store), который инициализируется при загрузке приложения через endpoint refresh.

\subsubsection{API-клиент}

Клиент для взаимодействия с серверным API реализует автоматическое обновление токена: при получении ответа 401 выполняется попытка обновления access token через \texttt{/auth/refresh}; если refresh-токен также недействителен, пользователь разлогинивается и перенаправляется на страницу входа.

Для повышения надёжности реализован механизм повторных попыток: до 3 попыток с экспоненциальной задержкой (500 мс базовая задержка + случайный джиттер) для ответов с кодом 5xx. API-клиент разделён на 23 модуля по доменным областям.

\subsubsection{Компоненты пользовательского интерфейса}

Библиотека переиспользуемых компонентов включает DataTable --- таблицу данных с сортировкой по столбцам, серверной пагинацией и адаптивным отображением на мобильных устройствах; Modal и ConfirmDialog --- модальные окна для создания, редактирования и подтверждения действий; FormField --- компонент формы с поддержкой валидации и отображения ошибок; FilterBar --- панель фильтров с поддержкой текстового поиска, выпадающих списков и диапазонов дат; MilkChart, LactationCurveChart, ReportChart, SeasonalChart --- компоненты визуализации данных на основе Chart.js; Toaster --- уведомления с раскрываемыми подробностями и автоматическим скрытием.

Валидация форм реализована на основе Svelte 5 Runes --- реактивной системы состояния, позволяющей декларативно описывать правила валидации и отслеживать состояние формы (dirty, errors, touched).

\subsubsection{Генерация отчётов}

Система поддерживает генерацию отчётов более 17 типов, покрывающих все аспекты управления стадом. Каждый отчёт доступен в формате JSON для отображения в интерфейсе, а также поддерживает экспорт в CSV и PDF. Ключевые типы отчётов перечислены в таблице~\ref{reports}.

\begin{table}
  \caption{Основные типы отчётов системы}
  \begin{tabular}{|c|l|p{6cm}|}
    \hline
    \textbf{Код} & \textbf{Название} & \textbf{Описание} \\
    \hline
    R12 & Здоровье вымени & Рабочий список коров с отклонениями за 24 часа \\
    \hline
    R13 & Неудачные доения & Визиты с неуспешным доением \\
    \hline
    R16 & Обзор стада & Дневная производительность стада \\
    \hline
    R18 & Остаток корма & Невыданный корм по коровам \\
    \hline
    R20 & Надой по времени & Суточный надой с разбивкой по показателям \\
    \hline
    R23 & Анализ вымени & Расширенный анализ за 14 дней \\
    \hline
    R24 & Активность/жвачка & Здоровье, активность и жвачка \\
    \hline
    R31-34 & Календарь & Ожидаемые отёлы, запуски, охота, стельность \\
    \hline
    R41 & Эффективность & Эффективность корова-робот \\
    \hline
    R52 & Лактация & Кривые лактации по номерам \\
    \hline
    R56 & Робот & Производительность роботов \\
    \hline
    R70 & Корм по типам & Стоимость и объём кормов по типам \\
    \hline
    R72 & Корм на корову & Потребление корма на корову \\
    \hline
  \end{tabular}
  \label{reports}
\end{table}
